\documentclass[11pt]{article}

\usepackage[margin=1in]{geometry}
\usepackage[T1]{fontenc}
\usepackage[utf8]{inputenc}
\usepackage{lmodern}
\usepackage{amsmath}
\usepackage{amssymb}
\usepackage{mathtools}
\usepackage{booktabs}
\usepackage{graphicx}
\usepackage{subcaption}
\usepackage{natbib}
\usepackage{hyperref}
\usepackage{siunitx}

\makeatletter
\IfFileExists{figures/sp500_train_val_test_split.pdf}{
  \newcommand{\paperroot}{}
}{
  \newcommand{\paperroot}{paper/}
}
\makeatother

\title{A Reproducible Comparison of Trend Estimation Methods on S\&P 500 Log Prices}
\author{Heriberto Espino-Montelongo}
\date{\today}

\begin{document}
\maketitle

\begin{abstract}
This paper compares a train-validation smoothness selector for penalized trend
estimation against classical and simple benchmark methods on historical S\&P
500 prices. The study uses the local \texttt{trend\_estimation} Python library
as a reproducible software dependency. Because real financial data do not
provide an observed true trend, the empirical comparison is based on temporal
out-of-sample forecasting errors and realized roughness diagnostics rather than
claims about true-trend recovery. The manuscript also distinguishes Guerrero
smoothness, an input index for finite-difference penalized smoothers, from
realized roughness, an output diagnostic available for every fitted trend.
\end{abstract}

\section{Introduction}
Trend estimation is a basic task in time series analysis, but trend smoothness
is rarely known in advance. Standard time-series and forecasting references
emphasize that trend extraction, prediction, and structural interpretation are
related but distinct goals \citep{chatfield2003analysis,brockwell2016introduction,
hamilton1994timeseries,makridakis1998forecasting}. This motivates procedures
that choose smoothness from data while preserving a strict temporal evaluation
protocol. This paper compares the proposed train-validation smoothness
selection approach with a time-weighted validation variant and several
classical or simple benchmarks.

The comparison is implemented as a reproducible subproject that uses
\texttt{trend\_estimation} as a normal Python library dependency
\citep{trendestimation}. The purpose is not to claim universal superiority from
one ticker, but to provide a concrete empirical case study and a template for
broader comparisons. The paper is deliberately explicit about terminology:
Guerrero smoothness is used only for finite-difference penalized smoothers,
whereas realized roughness is reported as a diagnostic for every fitted trend.

\section{Methods}
The proposed estimator uses penalized least squares. For an observed series
$y_1,\ldots,y_n$, the fitted trend $\hat{\tau}$ is obtained by balancing
fidelity to the observations against a finite-difference roughness penalty. A
typical objective has the form
\[
  \sum_{t=1}^n (y_t-\tau_t)^2 + \lambda \|D^d \tau\|_2^2,
\]
where $D^d$ is a difference operator of order $d$ and $\lambda$ controls
smoothness. This formulation follows the finite-difference penalized smoothing
tradition associated with Whittaker-Henderson graduation, Guerrero's
smoothness index, and modern penalized-spline and perfect-smoother treatments
\citep{whittaker1922graduation,henderson1924graduation,
guerrero2007penalized,eilers1996psplines,eilers2003perfect,wahba1990splines}.

The train-validation selector fits candidate penalized trends on the training
segment, forecasts the validation segment, and selects the candidate with the
lowest validation loss. The time-weighted variant uses a validation loss that
can emphasize later validation observations. The baseline methods are an
HP-style trend \citep{hodrick1997business,king1993lowfrequency}, a
Whittaker-Henderson-style penalized trend
\citep{whittaker1922graduation,henderson1924graduation}, a trailing moving
average, simple exponential smoothing, and a polynomial trend forecaster.
Moving-average and exponential-smoothing benchmarks are included because they
are common forecasting baselines \citep{brown1959forecasting,holt2004exponential,
winters1960forecasting,gardner1985exponential,gardner2006exponential,
hyndman2008exponential}. The polynomial benchmark is a low-dimensional
regression baseline \citep{draper1998regression}.

Guerrero's contribution is used here as a common smoothness scale for the
methods whose fitted trends solve a finite-difference penalized least-squares
problem. The train-validation selector and the time-weighted selector search
over Guerrero smoothness values and choose the value that minimizes their
validation loss. HPTrend and WhittakerTrend do not learn smoothness from the
validation set in this paper; instead, their conventional penalty choices are
converted back to the same Guerrero index so they can be interpreted on the
same model-side scale. Thus the four Guerrero-compatible methods share the
same penalty geometry, but they differ in how the penalty is selected.

The temporal protocol is strict. The series is split into 60 percent training,
20 percent validation, and 20 percent test observations. Hyperparameters are
selected using only the training and validation segments. Final test forecasts
are produced after refitting on the combined training and validation data.

\section{Algorithm}
The central algorithm fits a penalized trend for each candidate Guerrero
smoothness setting, evaluates out-of-sample validation forecasts, and then
refits the selected model before test evaluation. This section summarizes the
mechanism behind the penalized estimator, the numerical linear algebra used to
fit it, and the validation search used to choose among candidates. The
penalized objective belongs to the same finite-difference smoothing family as
Whittaker-Henderson graduation, Guerrero smoothing, and penalized spline
methods \citep{whittaker1922graduation,henderson1924graduation,
guerrero2007penalized,eilers1996psplines,eilers2003perfect,wahba1990splines}.

Let \(z=(z_1,\ldots,z_N)^\top\) denote the observed training series, let \(d\)
be the finite-difference order, and let
\[
  K \in \mathbb{R}^{(N-d)\times N}
\]
be the order-\(d\) differencing matrix. For a penalty value \(\lambda > 0\),
the penalized trend estimate is defined by
\begin{equation}
\label{eq:algorithm_pls}
  \hat{t}(\lambda)
  =
  \arg\min_{t\in\mathbb{R}^N}
  \left\{
    \|z-t\|_2^2 + \lambda \|K t\|_2^2
  \right\}.
\end{equation}
The first term measures fidelity to the observed series and the second term
penalizes roughness. Larger values of \(\lambda\) produce smoother fitted
trends.

The objective in \eqref{eq:algorithm_pls} is quadratic. Define
\[
  A(\lambda) = I_N + \lambda K^\top K.
\]
Taking the derivative with respect to \(t\) and setting it equal to zero gives
\[
  A(\lambda)\hat{t}(\lambda)=z.
\]
Since \(A(\lambda)\) is positive definite for \(\lambda>0\), the fitted trend is
\begin{equation}
\label{eq:algorithm_smoother}
  \hat{t}(\lambda)
  =
  S(\lambda)z,
  \qquad
  S(\lambda)=\left(I_N+\lambda K^\top K\right)^{-1}.
\end{equation}
Here \(S(\lambda)\) is the smoothing matrix.

\subsection{Spectral solution of the penalized system}
Although \eqref{eq:algorithm_smoother} is written with an inverse, the
implementation does not form a dense inverse directly. For a fixed sample size
\(N\) and order \(d\), it first constructs the penalty matrix
\[
  P_d = K^\top K.
\]
This matrix is symmetric positive semidefinite, so it admits an orthogonal
eigendecomposition
\[
  P_d = Q \Delta Q^\top,
  \qquad
  \Delta=\mathrm{diag}(\delta_1,\ldots,\delta_N).
\]
Then
\[
  I_N+\lambda P_d
  =
  Q(I_N+\lambda \Delta)Q^\top,
\]
and the fitted trend can be computed as
\begin{equation}
\label{eq:algorithm_spectral_solution}
  \hat{t}(\lambda)
  =
  Q
  \mathrm{diag}\left(
    \frac{1}{1+\lambda \delta_1},
    \ldots,
    \frac{1}{1+\lambda \delta_N}
  \right)
  Q^\top z.
\end{equation}
This reduces each candidate fit to projections onto the eigenvectors, elementwise
division by \(1+\lambda\delta_i\), and a projection back to the original time
domain. Since the same \(Q\) and \(\delta_i\) are reused for many candidate
\(\lambda\) values at the same \(N\) and \(d\), the paper pipeline benefits from
caching the solver for repeated fits. The same eigenvalues also provide the
effective-degrees-of-freedom trace used in Guerrero's index, so the numerical
linear algebra for fitting and the numerical linear algebra for tuning are
shared rather than duplicated.

For Guerrero-style penalized trends, the implementation may also estimate a
constant \(d\)-th difference drift. In that case the right-hand side is adjusted
iteratively and the drift estimate \(m\) is updated from the fitted trend until
the change is smaller than a numerical tolerance. HPTrend and WhittakerTrend are
run as classical fixed-drift baselines with this drift set to zero.

\subsection{Mapping Guerrero smoothness to \texorpdfstring{\(\lambda\)}{lambda}}
The paper reports Guerrero smoothness for finite-difference penalized smoothers,
but the linear system is solved in terms of \(\lambda\). The eigendecomposition
above gives a numerically convenient expression for the trace of the smoothing
matrix:
\[
  \mathrm{tr}\{S(\lambda)\}
  =
  \sum_{i=1}^N \frac{1}{1+\lambda\delta_i}.
\]
For \(d\ge 1\), the implementation uses the normalized Guerrero index
\begin{equation}
\label{eq:algorithm_guerrero_normalized}
  s(\lambda;N,d)
  =
  \frac{
    1 - N^{-1}\sum_{i=1}^N (1+\lambda\delta_i)^{-1}
  }{
    1-d/N
  },
\end{equation}
which maps the attainable range to a scale close to \([0,1)\) for the finite
sample and difference order \citep{guerrero2007penalized}. Given a requested
Guerrero smoothness value \(s\), the code finds the corresponding \(\lambda\)
by bisection:
\begin{enumerate}
  \item Set an initial bracket \([0,1]\).
  \item Increase the upper endpoint by factors of ten until the target
        smoothness is bracketed.
  \item Bisect the interval until the trace-based smoothness is within the
        numerical tolerance.
\end{enumerate}
This step is monotone because increasing \(\lambda\) decreases the effective
degrees of freedom \(\mathrm{tr}\{S(\lambda)\}\) and therefore increases the
Guerrero smoothness index.

\subsection{How Guerrero's idea defines the penalized methods}
Guerrero's idea is not a separate estimator in the pipeline; it is the
trace-based scale used to parameterize and compare finite-difference penalized
smoothers. Once \(N\) and \(d\) are fixed, every penalty value \(\lambda\)
corresponds to a Guerrero smoothness value through the smoothing matrix trace.
The implementation uses this correspondence in four related ways.

For TrainValidationSelector, the candidate set is written directly as pairs
\((d,s)\), where \(s\) is the Guerrero smoothness index. Each candidate \(s\)
is converted to its corresponding \(\lambda\), the penalized system is solved
on the training segment, and the resulting trend is extrapolated into the
validation segment. The selected model is the candidate with the smallest
validation error. In this sense, Guerrero's index becomes the hyperparameter
that the temporal validation split chooses.

For TimeWeightedValidationSelector, the same Guerrero-indexed candidate set is
used. The difference is only in the validation criterion: later validation
observations receive larger weights, so the selected Guerrero smoothness is the
one that performs best for the more recent part of the validation window. The
underlying finite-difference penalty, the \(s\)-to-\(\lambda\) conversion, and
the linear solve are otherwise the same as in TrainValidationSelector.

For HPTrend, the trend is also a finite-difference penalized least-squares
smoother, usually with a second-difference penalty. The classical HP
presentation specifies a penalty \(\lambda\) rather than a Guerrero smoothness
index. In the paper pipeline, that fixed \(\lambda\) is passed through the same
trace formula to report the implied Guerrero smoothness. HPTrend is therefore
Guerrero-compatible for interpretation, even though the validation set is not
used to choose its penalty.

For WhittakerTrend, the fitted trend again solves the same algebraic family of
linear systems, with a finite-difference penalty and a fixed penalty setting.
As with HPTrend, the fixed \(\lambda\) determines a fitted smoother matrix and
therefore an implied Guerrero smoothness index. This lets WhittakerTrend be
reported beside the validation-selected methods without pretending that it ran
a validation search.

This is why the paper reports Guerrero smoothness for TrainValidationSelector,
TimeWeightedValidationSelector, HPTrend, and WhittakerTrend, but leaves that
field blank for moving averages, exponential smoothing, and polynomial
forecasting. Those latter methods can still be evaluated by forecast error and
realized roughness, but they do not have a native \(S_d(\lambda;N)\) because
they are not defined by the penalized system in
\eqref{eq:algorithm_pls}.

For validation, let \(R\) be the matrix that selects validation observations and
let \(n_{\mathrm{val}}\) be the number of validation points. The validation
residual is
\[
  r(\lambda) = Rz - R S(\lambda)z,
\]
and the validation mean squared error is
\begin{equation}
\label{eq:algorithm_validation_mse}
  f(\lambda)
  =
  \frac{1}{n_{\mathrm{val}}}
  r(\lambda)^\top r(\lambda).
\end{equation}
The derivative calculation in the companion report follows from
\[
  S'(\lambda)
  =
  -S(\lambda)K^\top K S(\lambda),
\]
which implies
\[
  r'(\lambda)
  =
  R S(\lambda)K^\top K S(\lambda)z.
\]
Therefore,
\begin{equation}
\label{eq:algorithm_validation_derivative}
  f'(\lambda)
  =
  \frac{2}{n_{\mathrm{val}}}
  \left[ R\left(I_N-S(\lambda)\right)z \right]^\top
  \left[ R S(\lambda)K^\top K S(\lambda)z \right].
\end{equation}

\subsection{Validation search and local minima}
The derivative in \eqref{eq:algorithm_validation_derivative} explains the
shape of the validation objective, but the implemented paper pipeline uses a
robust direct search over candidate Guerrero smoothness values. For each
candidate difference order \(d\) and smoothness value \(s\), the algorithm:
\begin{enumerate}
  \item converts \(s\) to \(\lambda\) using the bisection procedure above;
  \item fits the penalized trend on the training segment only;
  \item extrapolates the fitted trend into the validation segment;
  \item computes the validation loss between the forecast and the observed
        validation log prices.
\end{enumerate}
For the S\&P 500 comparison, the paper script evaluates orders
\(d\in\{1,2,3\}\) and a fixed grid of Guerrero smoothness values between 0.05
and 0.98. The broader library default uses a denser grid, but the paper uses a
smaller reproducible grid to keep the empirical study fast and stable.

The selector records the full validation curve and detects local minima on the
one-dimensional smoothness grid for each order. A grid point is treated as a
local minimum when its validation score is no larger than the scores of its
neighbors; boundary points can also be local minima if they improve on their
single neighbor. The library can optionally refine each detected bracket by
golden-section search. In the paper pipeline, refinement is disabled so that the
saved validation curves, local minima, and selected value all correspond exactly
to the stated reproducible grid.

The time-weighted selector uses the same candidate search, but replaces the
ordinary mean squared validation loss with a weighted validation loss. In the
paper run, the weights follow an exponential decay scheme, so later validation
observations receive relatively more emphasis than earlier validation
observations. This changes only the scoring rule, not the penalized smoothing
linear algebra. In forecasting terminology, this is a validation-rule change
rather than a different smoothing model \citep{makridakis1998forecasting,
hyndman2008exponential}.

\subsection{Forecasting from a fitted penalized trend}
After fitting a penalized trend on the training segment, the package forecasts
future trend values by extending the finite-difference structure at the right
tail. If the selected order is \(d\), the forecast recursively imposes a
constant estimated \(d\)-th difference \(m\). For \(d=1\), this is a constant
slope extension; for \(d=2\), it is a constant second-difference extension; and
for higher orders it extends the corresponding finite-difference recurrence.
These forecasts are used for validation scoring and, after refitting on
training plus validation, for final test evaluation. The benchmark forecasts
are intentionally simpler: trailing moving averages, exponential smoothing, and
polynomial extrapolation are used as transparent baselines rather than as a
complete forecasting taxonomy \citep{brown1959forecasting,holt2004exponential,
winters1960forecasting,gardner1985exponential,gardner2006exponential,
draper1998regression}.

\subsection{Final comparison protocol}
The complete numerical procedure for selector-based penalized methods is:
\begin{enumerate}
  \item split the series temporally into training, validation, and test
        segments;
  \item fit and score each \((d,s)\) candidate on training-to-validation
        forecasts;
  \item choose the candidate with the lowest validation loss, without using the
        test segment;
  \item refit the chosen penalized smoother on training plus validation;
  \item forecast the test segment and compute MAE, RMSE, SMAPE, and related
        metrics;
  \item compute realized roughness \(R_d(\hat{t})=\|D^d\hat{t}\|_2^2\) from the
        fitted training-plus-validation trend.
\end{enumerate}
The non-Guerrero baselines use the same temporal protocol, but they do not have
a Guerrero smoothness search. The moving-average baseline uses its fixed
trailing window, the exponential-smoothing baseline uses its fixed \(\alpha\),
and the polynomial forecaster uses its fixed degree. They still receive test
forecasts and realized roughness diagnostics, but their Guerrero smoothness
entries are left blank.

The Guerrero smoothness index used by the implementation is
\[
  S_d(\lambda; N)
  =
  1 -
  \frac{\mathrm{tr}\left[S(\lambda)\right]}{N},
  \qquad
  S(\lambda)=\left(I_N+\lambda K^\top K\right)^{-1}.
\]
This index is meaningful for the finite-difference penalized smoothers in the
comparison: the train-validation selector, the time-weighted selector, HPTrend,
and WhittakerTrend. It is not a native parameter of moving averages,
exponential smoothing, or polynomial forecasting.

In practice, solving \(f'(\lambda)=0\) does not give a simple general closed
form. The implementation therefore uses numerical validation search over the
package's normalized Guerrero smoothness parameter, converts the selected index
to \(\lambda\), and then refits the final penalized trend on the combined
training and validation sample. The comparative pipeline applies the same
temporal selection principle to the proposed train-validation selector and its
time-weighted variant, while fixed-parameter baselines are fit on the same
training and refit samples.

\section{Data}
The empirical data are daily S\&P 500 observations from \texttt{yfinance}
\citep{yfinance}, using ticker \texttt{\string^GSPC} from 2010-01-01 onward.
The analysis uses
\[
  y_t = \log(P_t),
\]
where $P_t$ is the adjusted close price when available and the close price
otherwise. The raw and processed data are cached under \texttt{paper/data/} so
that the experiment does not require a download on every run.

\section{Metrics}
Forecast accuracy is reported using MAE, MSE, RMSE, MAPE, and SMAPE. The main
target is the observed log-price series. Real financial data do not provide an
observed true trend, so this report does not evaluate true-trend recovery on the
S\&P 500 series. These forecast-error summaries are standard descriptive tools
in empirical forecasting comparisons \citep{makridakis1998forecasting,
hyndman2008exponential}.

For finite-difference penalized least-squares smoothers, the reported
Guerrero smoothness value is the model-side index
\[
  S_d(\lambda; N)
  =
  1 -
  \frac{\mathrm{tr}\left[\left(I_N+\lambda D_d^\top D_d\right)^{-1}\right]}{N},
  \qquad d \ge 1.
\]
It is used to choose or interpret the penalty \(\lambda\) once the sample size
\(N\) and difference order \(d\) are fixed \citep{guerrero2007penalized}.

As a separate output diagnostic, the paper reports realized roughness
\[
  R_d(\hat{\tau}) = \|D^d\hat{\tau}\|_2^2.
\]
Realized roughness can be computed for any fitted trend and should not be
identified with Guerrero smoothness.

\section{Results}
Figure~\ref{fig:split} shows the temporal split used in the experiment.
Figure~\ref{fig:trends} shows the fitted trends before the test period.
Figure~\ref{fig:validation} shows validation curves for the selector-based
methods. Figure~\ref{fig:forecasts} displays test forecasts, while
Figures~\ref{fig:rmse} and~\ref{fig:smape} summarize test errors. The realized
roughness diagnostic is shown in Figure~\ref{fig:roughness}.

\begin{figure}[htbp]
  \centering
  \includegraphics[width=\linewidth]{\paperroot figures/sp500_train_val_test_split.pdf}
  \caption{Temporal train, validation, and test split for S\&P 500 log prices.}
  \label{fig:split}
\end{figure}

Figure~\ref{fig:split} documents the chronological evaluation design. The
training segment is used to fit candidate models, the validation segment is
used only for hyperparameter selection, and the test segment is held out until
final evaluation. This ordering is important because randomly mixing financial
time points would leak future information into the model-selection step.

\begin{figure}[htbp]
  \centering
  \includegraphics[width=\linewidth]{\paperroot figures/sp500_smoothed_trends.pdf}
  \caption{In-sample trend estimates fitted before the test period.}
  \label{fig:trends}
\end{figure}

Figure~\ref{fig:trends} compares the fitted trends on the train-plus-validation
sample before test forecasting begins. Curves that stay close to the observed
series have lower in-sample smoothing bias but may carry more short-run
variation, while flatter curves impose stronger smoothness. The figure is
therefore a visual complement to the parameter and realized-roughness tables:
it shows what each smoothing choice does to the observed log-price path.

\begin{figure}[htbp]
  \centering
  \includegraphics[width=\linewidth]{\paperroot figures/sp500_validation_curves.pdf}
  \caption{Validation curves over the Guerrero smoothness index for selector-based penalized methods. The vertical scale is clipped to keep the lower-order curves readable when the order-three curve is much larger.}
  \label{fig:validation}
\end{figure}

Figure~\ref{fig:validation} shows how validation error changes as the Guerrero
smoothness index varies for each candidate difference order. The selected point
is the minimum of the displayed validation criterion for each selector. In this
run, the order-three curves are much larger than the order-one and order-two
curves, so the plotted y-axis is limited to reveal the smaller-scale structure
that determines the useful comparison among the better-performing candidates.

\begin{figure}[htbp]
  \centering
  \includegraphics[width=\linewidth]{\paperroot figures/sp500_forecasts_test.pdf}
  \caption{Test-set forecasts from the compared methods. Pale line segments show each method's fitted train-validation trend with alpha 0.2; bright segments show the held-out test forecasts.}
  \label{fig:forecasts}
\end{figure}

Figure~\ref{fig:forecasts} overlays each method's test-period forecast on the
held-out observed log prices. The pale pre-test segments show the fitted trend
state from which each method extrapolates, which explains why the bright
forecast segments do not necessarily begin at the last raw observed value. This
figure should be read together with the test-error table: a forecast that
visually tracks the level and slope of the test period tends to have lower RMSE
and SMAPE, while forecasts that remain too flat or extrapolate the wrong slope
accumulate larger errors.

\begin{figure}[htbp]
  \centering
  \begin{subfigure}{0.48\linewidth}
    \includegraphics[width=\linewidth]{\paperroot figures/sp500_test_rmse_barplot.pdf}
    \caption{RMSE}
    \label{fig:rmse}
  \end{subfigure}
  \begin{subfigure}{0.48\linewidth}
    \includegraphics[width=\linewidth]{\paperroot figures/sp500_test_smape_barplot.pdf}
    \caption{SMAPE}
    \label{fig:smape}
  \end{subfigure}
  \caption{Test error summaries.}
\end{figure}

Figures~\ref{fig:rmse} and~\ref{fig:smape} rank methods by two test-set error
summaries. RMSE emphasizes larger forecast mistakes because errors are squared
before averaging, while SMAPE reports a scale-free percentage-style error. A
method that performs well in both panels is more robustly accurate on the
chosen test window than a method that improves only one metric.

\begin{figure}[htbp]
  \centering
  \includegraphics[width=0.8\linewidth]{\paperroot figures/sp500_roughness_vs_error.pdf}
  \caption{Realized roughness diagnostic against test RMSE.}
  \label{fig:roughness}
\end{figure}

Figure~\ref{fig:roughness} separates forecast performance from realized
roughness. The horizontal axis is the diagnostic
\(R_d(\hat{\tau})=\|D^d\hat{\tau}\|_2^2\), computed after fitting each trend;
it is not Guerrero smoothness. Points far to the left are smoother in realized
finite-difference terms, while lower points have better test RMSE. The scatter
therefore shows whether extra smoothness coincides with better forecasting in
this empirical split.

\subsection{Tables}
Table~\ref{tab:test_metrics} reports the final held-out test results. The
columns MAE, RMSE, and SMAPE summarize forecast accuracy on observations that
were not used for hyperparameter selection. The parameter columns should be
read method by method: Guerrero smoothness and \(\lambda\) are meaningful only
for finite-difference penalized smoothers, while window, alpha, and degree
describe the moving-average, exponential-smoothing, and polynomial baselines.

\begin{table}[htbp]
  \centering
  \caption{Held-out test accuracy and method-specific settings.}
  \label{tab:test_metrics}
  \begin{center}
\resizebox{\linewidth}{!}{%
\begin{tabular}{lllllllllll}
\toprule
model & n\_obs & MAE & RMSE & SMAPE & best\_order & guerrero\_smoothness & best\_lambda & window & alpha & degree \\
\midrule
TrainValidationSelector & 827 & 0.127263 & 0.148216 & 1.475928 & 1.0 & 0.05 & 0.027008 &  &  &  \\
TimeWeightedValidationSelector & 827 & 0.147868 & 0.170415 & 1.717913 & 1.0 & 0.98 & 615.756247 &  &  &  \\
HPTrend & 827 & 0.675167 & 0.774414 & 7.389961 & 2.0 & 0.944195 & 1600.0 &  &  &  \\
WhittakerTrend & 827 & 0.039826 & 0.04672 & 0.462446 & 2.0 & 0.75 & 4.874751 &  &  &  \\
MovingAverageTrend & 827 & 0.322219 & 0.367195 & 3.792746 &  &  &  & 60.0 &  &  \\
ExponentialSmoothingTrend & 827 & 0.290947 & 0.338309 & 3.416307 &  &  &  &  & 0.2 &  \\
PolynomialTrendForecaster & 827 & 0.083899 & 0.09498 & 0.973113 &  &  &  &  &  & 2.0 \\
\bottomrule
\end{tabular}%
}
\end{center}

\end{table}

Table~\ref{tab:selection_metrics} reports validation-period performance before
the final test refit. These values explain why the validation-based selectors
chose their reported settings, but they are not final generalization errors.
Comparing this table with Table~\ref{tab:test_metrics} shows whether a method
that looked good during validation continued to perform well on the held-out
test window.

\begin{table}[htbp]
  \centering
  \caption{Validation-period accuracy used for model-selection comparison.}
  \label{tab:selection_metrics}
  \begin{center}
\resizebox{\linewidth}{!}{%
\begin{tabular}{lllllllllll}
\toprule
model & n\_obs & MAE & RMSE & SMAPE & best\_order & guerrero\_smoothness & best\_lambda & window & alpha & degree \\
\midrule
TrainValidationSelector & 825 & 0.096959 & 0.118151 & 1.176689 & 1.0 & 0.05 & 0.027008 &  &  &  \\
TimeWeightedValidationSelector & 825 & 0.103816 & 0.126905 & 1.260935 & 1.0 & 0.98 & 615.756247 &  &  &  \\
HPTrend & 825 & 0.627408 & 0.742694 & 7.196558 & 2.0 & 0.944195 & 1600.0 &  &  &  \\
WhittakerTrend & 825 & 0.939478 & 1.099213 & 10.529626 & 2.0 & 0.75 & 4.874751 &  &  &  \\
MovingAverageTrend & 825 & 0.25886 & 0.290361 & 3.17557 &  &  &  & 60.0 &  &  \\
ExponentialSmoothingTrend & 825 & 0.237732 & 0.269659 & 2.911698 &  &  &  &  & 0.2 &  \\
PolynomialTrendForecaster & 825 & 0.133972 & 0.161477 & 1.629947 &  &  &  &  &  & 2.0 \\
\bottomrule
\end{tabular}%
}
\end{center}

\end{table}

Table~\ref{tab:model_parameters} collects the selected or fixed model settings
in one place. The selector rows show the chosen difference order, Guerrero
smoothness index, and corresponding \(\lambda\). HPTrend and WhittakerTrend
show their fixed finite-difference penalty settings and implied Guerrero
indices. The non-Guerrero benchmarks deliberately have blank Guerrero
smoothness entries, because their tuning parameters are window length,
exponential alpha, or polynomial degree rather than \(S_d(\lambda;N)\).

\begin{table}[htbp]
  \centering
  \caption{Selected or fixed model parameters.}
  \label{tab:model_parameters}
  \begin{center}
\resizebox{\linewidth}{!}{%
\begin{tabular}{llllllll}
\toprule
model & kind & best\_order & guerrero\_smoothness & best\_lambda & degree & window & alpha \\
\midrule
TrainValidationSelector & selector & 1.0 & 0.05 & 0.027008 &  &  &  \\
TimeWeightedValidationSelector & selector & 1.0 & 0.98 & 615.756247 &  &  &  \\
HPTrend & estimator & 2.0 & 0.944195 & 1600.0 &  &  &  \\
WhittakerTrend & estimator & 2.0 & 0.75 & 4.874751 &  &  &  \\
MovingAverageTrend & estimator &  &  &  &  & 60.0 &  \\
ExponentialSmoothingTrend & estimator &  &  &  &  &  & 0.2 \\
PolynomialTrendForecaster & polynomial\_forecaster &  &  &  & 2.0 &  &  \\
\bottomrule
\end{tabular}%
}
\end{center}

\end{table}

Table~\ref{tab:roughness_metrics} reports realized roughness for every fitted
trend using the stated roughness order. Unlike Guerrero smoothness, this
diagnostic is computed from the fitted trend itself and can therefore be
reported for all methods. The accompanying test RMSE column makes it possible
to compare the smoothness of the fitted curve with its forecasting accuracy.

\begin{table}[htbp]
  \centering
  \caption{Realized roughness diagnostics and test RMSE.}
  \label{tab:roughness_metrics}
  \begin{center}
\resizebox{\linewidth}{!}{%
\begin{tabular}{llll}
\toprule
model & roughness\_order & realized\_roughness & test\_RMSE \\
\midrule
TrainValidationSelector & 1 & 0.372955 & 0.148216 \\
TimeWeightedValidationSelector & 1 & 0.002511 & 0.170415 \\
HPTrend & 2 & 0.000126 & 0.774414 \\
WhittakerTrend & 2 & 0.010925 & 0.04672 \\
MovingAverageTrend & 2 & 0.000254 & 0.367195 \\
ExponentialSmoothingTrend & 2 & 0.018961 & 0.338309 \\
PolynomialTrendForecaster & 2 & 0.0 & 0.09498 \\
\bottomrule
\end{tabular}%
}
\end{center}

\end{table}

\section{Discussion}
The results should be interpreted as one empirical comparison on one financial
index over one historical period. Forecasting accuracy and trend recovery are
different goals, and a method that forecasts well over the selected test window
need not recover a latent economic trend. Validation-based selection can also
overfit the validation period, especially when the candidate grid is wide or
the market regime changes between validation and test.

Guerrero smoothness is a model-side index defined only for finite-difference
penalized least-squares smoothers. It maps a smoothing penalty \(\lambda\) into
a comparable scale \(S_d(\lambda; N)\) once the sample size \(N\) and
difference order \(d\) are fixed. By contrast, realized roughness is an output
diagnostic \(R_d(\hat{\tau}) = \|D^d\hat{\tau}\|_2^2\) that can be computed for
any fitted trend. Therefore, moving-average, exponential-smoothing, and
polynomial benchmarks do not have a native Guerrero smoothness index, although
their fitted trends can still be compared through realized roughness.

Future sensitivity checks should vary the time period, include additional
tickers, and compare against synthetic experiments where the true trend is
known. Additional baselines such as Kalman filters, structural time-series
models, ARIMA models, STL/LOESS, smoothing splines, and structural-break models
would also make the comparison broader
\citep{kalman1960new,harvey1989structural,box2008time,cleveland1979lowess,
cleveland1990stl,wahba1990splines}.

\section{Reproducibility}
The empirical pipeline is distributed with the local \texttt{trend\_estimation}
package and paper scripts \citep{trendestimation}. From the repository root,
activate the environment and install the local package:
\begin{verbatim}
conda activate trend_estimation
python -m pip install -e ".[dev,finance]"
\end{verbatim}
Then run:
\begin{verbatim}
python paper/build_paper.py --ticker ^GSPC --start 2010-01-01 --no-compile
\end{verbatim}
If \texttt{latexmk} is installed, compile with:
\begin{verbatim}
latexmk -pdf -interaction=nonstopmode -outdir=paper/build paper/main.tex
\end{verbatim}

\section{Conclusion}
This paper establishes a reproducible comparison between train-validation
penalized trend selection, a time-weighted variant, and several benchmark
methods on S\&P 500 log prices. The current evidence is descriptive and
case-specific. Broader empirical and synthetic experiments are needed before
making stronger methodological claims.


\bibliographystyle{plainnat}
\IfFileExists{biblio.bib}{
  \bibliography{biblio}
}{
  \bibliography{paper/biblio}
}

\end{document}
